% !TeX program = xelatex
% ============================================================================
%  第 11 课　幂的尾巴转圈了：从现象到猜想　—— 教师版讲义（原 14+15 合并）
%  与 02-课件/第11课-幂的尾巴转圈了/第11课-幂的尾巴转圈了.tex 同步
% ============================================================================

\jhlesson{11}{幂的尾巴转圈了：从现象到猜想}{时长 45 分钟\quad·\quad 教具：无}

\begin{mubiao}
  \begin{itemize}
    \item 会算 \(a\) 的幂除以 \(p\) 的余数，并发现它在循环。
    \item 能在好几个例子里发现同一个现象，自己归纳出猜想。
  \end{itemize}
\end{mubiao}

\noindent{\small 本课节奏：0—6 回顾｜6—16 算 \(2\) 的幂｜16—30 换几个数试试｜30—38 归纳猜想｜38—45 收尾。}

\section*{一、回顾（0—6 分钟）}

短短几句把「绕圈」的三次见面串起来：钟表走满 \(12\) 格回到原地；
三角形、正方形的动作做几次回到原位；圆盘上一直加同一个数，\(\langle3\rangle\) 只圈住 \(4\) 个。

一句话过渡：今天看第四种绕圈——\textbf{幂的尾巴}。

\section*{二、算 \(2\) 的幂（6—16 分钟）}

板书两行，第二行现场算：

\begin{center}
\begin{tabular}{@{}lcccccc@{}}
  \toprule
  \(n\) & \(1\) & \(2\) & \(3\) & \(4\) & \(5\) & \(6\) \\
  \midrule
  \(2^n\) & \(2\) & \(4\) & \(8\) & \(16\) & \(32\) & \(64\) \\
  除以 \(7\) 的余数 & \(2\) & \(4\) & \(1\) & \(2\) & \(4\) & \(1\) \\
  \bottomrule
\end{tabular}
\end{center}

\begin{caozuo}
  余数一行只写前两个（\(2,4\)），然后让全班喊下一个。喊出 \(1\) 之后继续喊，
  他们很快会自己发现：\textbf{从第 \(4\) 个开始原样重复}。
\end{caozuo}

跟上节课那句对一下：这和圆盘上一直加是一回事，只是把「加」换成了「乘」。
这里乘法的主角是 \(2\)，它生成的圈子是 \(2,4,1\)。

\begin{zhu}
  学生版「想一想」两问的答案口径：第 \(7\) 个余数是 \(2\)
  （\(2^{7}=128=18\times7+2\)，和第 \(1\) 个一样）；第 \(2026\) 个也是 \(2\)
  （\(2026=3\times675+1\)，落在循环 \(2,4,1\) 的第 \(1\) 位）。
  先带他们做第 \(7\) 个，让他们自己发现「只要看指数除以 \(3\) 的余数」，再问第 \(2026\) 个。
\end{zhu}

\section*{三、换几个数试试（16—30 分钟）}

\begin{caozuo}
  \textbf{全班在自己讲义上算几组。}每组只算到「第一次出现 \(1\)」为止，不必算完。
  算完请他们用手指比出「第几步到 \(1\)」，你往黑板上的总表里填。
\end{caozuo}

\begin{center}
\begin{tabular}{@{}ccc@{}}
  \toprule
  \(p\) & \(a\) & 余数序列（到第一次出现 \(1\) 为止） \\
  \midrule
  \(5\) & \(2\) & \(2,\ 4,\ 3,\ 1\) \\
  \(5\) & \(3\) & \(3,\ 4,\ 2,\ 1\) \\
  \(7\) & \(2\) & \(2,\ 4,\ 1\) \\
  \(7\) & \(3\) & \(3,\ 2,\ 6,\ 4,\ 5,\ 1\) \\
  \bottomrule
\end{tabular}
\end{center}

\begin{zhu}
  算 \(3^n\) 除以 \(7\) 的时候，学生会嫌数大。提醒他们\emph{每次只用上一步的余数去乘}：
  \(3\to3\times3=9\equiv2\to2\times3=6\to6\times3=18\equiv4\to\dots\)
  这样数字永远不超过 \(7\)，一点也不难算。
\end{zhu}

\section*{四、归纳猜想（30—38 分钟）}

把「第几步回到 \(1\)」单独列一张表（按 \(p\) 分组讲：\(p=5\) 两组、\(p=7\) 四组、\(p=11\) 两组 + \(p=13\) 两组）：

\begin{center}
\begin{tabular}{@{}ccc@{}}
  \toprule
  \(p\) & \(a\) & 第几步第一次回到 \(1\) \\
  \midrule
  \(5\) & \(2\) & \(4\) \\
  \(5\) & \(3\) & \(4\) \\
  \(7\) & \(2\) & \(3\) \\
  \(7\) & \(3\) & \(6\) \\
  \(7\) & \(4\) & \(3\) \\
  \(7\) & \(5\) & \(6\) \\
  \(11\) & \(2\) & \(10\) \\
  \(11\) & \(3\) & \(5\) \\
  \(13\) & \(2\) & \(12\) \\
  \(13\) & \(3\) & \(3\) \\
  \bottomrule
\end{tabular}
\end{center}

表里前四行（\(p=5\) 的两组、\(p=7\) 的前两组）是全班自己算的，后六行是你补的——
说「\(10\) 组」的时候要分清这两块（学生只亲手算了 \(4\) 组），学生版那张表的说明与此一致。

\begin{caozuo}
  不要直接给答案，让全班自己看两分钟，然后请人起来说。引导到两句话：
  \begin{itemize}
    \item 回到 \(1\) 的步数，都\textbf{整除} \(p-1\)（\(p=7\) 时整除 \(6\)，\(p=5\) 时整除 \(4\)）；
    \item 那么走满 \(p-1\) 步时，一定是 \(1\)：\(2^6\equiv1\)、\(3^6\equiv1\)、\(2^4\equiv1\)、\(3^4\equiv1\)。
  \end{itemize}
\end{caozuo}

\begin{dingli}[猜想]
  设 \(p\) 是\textbf{质数}（也叫素数——就是只能被 \(1\) 和它自己整除的数），\(a\) 不是 \(p\) 的倍数，那么
  \[
    a^{\,p-1}\equiv1 \pmod p .
  \]
\end{dingli}

\begin{zhu}
  强调一句：这些例子都对，但「试出来的」还不算数——必须证明。
  这句话为下一节课留出悬念。
\end{zhu}

\section*{五、收尾（38—45 分钟）}

板书三句话：幂的余数会绕圈，且一定回到 \(1\)；回到 \(1\) 的步数整除 \(p-1\)；
由此猜出 \(a^{p-1}\equiv1\pmod p\)。

\begin{xiaojie}
  \begin{itemize}
    \item 幂的余数会绕圈：\(2^n\) 除以 \(7\) 是 \(2,4,1,2,4,1,\dots\)。
    \item 回到 \(1\) 的步数整除 \(p-1\)。
    \item 猜想：\(p\) 是质数、\(a\) 不是 \(p\) 的倍数时，\(a^{p-1}\equiv1\pmod p\)。
  \end{itemize}
\end{xiaojie}

\noindent\textbf{下节课}：把「猜」换成「证」——费马小定理的证明。